\section{A Conditional Model of Decay}
\label{sec:decay_model}
\subsection{Latency, drag, and capacity}
For a protocol $s$ with structural cost $K(s)$ and update frequency $f$, adopt
\begin{equation}
\Phi_{\rm drag}=\xi K(s)f.
\end{equation}
The capacity compared with this drag must have the same units. A Landauer bound concerns logically irreversible erasure under specified thermodynamic conditions; it is not a universal cost for every reversible update. Interpreting excess demand as a collision of lattice writes requires a scheduling and encoding model, and does not itself determine a decay rate.

\subsection{Allowed transitions and stability}
Let $L(s)\in\mathbb N$ be latency, $q(s)$ a conserved label, and $s\to t$ an allowed transition. Assume
\begin{equation}
s\to t\ \Longrightarrow\ L(t)<L(s),\qquad q(t)=q(s).
\end{equation}
The label describes the entire exit cohort, not necessarily each daughter. A terminal protocol has no outgoing transition. A latency minimum within its label class is terminal, and a sequence of strict reductions has at most $L(s_0)$ steps. A lower-latency configuration need not be reachable: amplitudes, selection rules and thresholds are additional conditions. These routing assumptions alone do not prove absolute electron or proton stability.

\subsection{Widths, rates and waiting times}
Write $\Gamma$ for an energy width and $\gamma=\Gamma/\hbar$ for a rate. For $\Gamma>0$,
\begin{equation}
\tau=\frac{1}{\gamma}=\frac{\hbar}{\Gamma}.
\end{equation}
A stationary, memoryless hazard gives survival
\begin{equation}
S(t)=e^{-\gamma t},\quad S(t+u)=S(t)S(u),\quad
 t_{1/2}=\frac{\log 2}{\gamma}.
\end{equation}
Poisson event counts have exponential waiting times; the waiting time itself is not Poisson-distributed. Deterministic routing with unspecified background traffic does not establish this law. Deriving a stationary hazard from microscopic dynamics remains an open task.
